% ─────────────────────────────────────────────────────────────────────────────
% Section 1: Introduction
% Paper: "Silent Decisions Are Type Errors: Enforcing AI Accountability
%         via Linear Resource Types"
% Venue: IEEE Computer — Special Issue on AI Governance
% Deadline: 13 April 2026
% Authors: Daniel Lau Pereira Soares
% ─────────────────────────────────────────────────────────────────────────────

\section{Introduction}

Deployed AI systems now issue consequential decisions in milliseconds:
a credit application is rejected, a job candidate is ranked out of contention,
a social-benefit disbursement is blocked.
Regulatory frameworks acknowledge this reality.
Brazil's Lei Geral de Proteção de Dados (LGPD), Article~18, \S2,
grants affected individuals the right to demand human review and explanation
of any automated decision~\cite{brazil2018lgpd}.
The European Union AI Act, Article~86, extends that obligation to high-risk
AI systems across the bloc~\cite{eu2024aiact}.
Yet both instruments share a structural gap: they mandate that explanations
and contestability rights \emph{accompany} a decision,
but they cannot prevent a decision from existing \emph{without} them.
A system may emit a denial, log it asynchronously, and silently drop the log.
No law is formally broken in that instant; the violation surfaces only later,
if at all.
We call this phenomenon a \emph{silent decision}:
an AI-generated act with legal consequences for an individual
that carries no non-repudiable evidence chain.
Silent decisions are not a failure of intent---they are a failure of type.

The conventional response is instrumentation.
Policy engines such as Open Policy Agent~\cite{opa2024} and
Kyverno govern cloud infrastructure---pods, secrets, network policies---issuing
allow/deny verdicts over \emph{resources}, not over \emph{persons}.
Evidence of a decision's basis is a separable artifact, appended
asynchronously, omissible under load, and falsifiable after the fact.
Runtime assertions offer a second line of defense:
\texttt{assert(evidence != null)} at every callsite.
The analogy to memory safety is instructive.
In~C, \texttt{free()} can be paired with discipline and checked by analyzers---
yet dangling-pointer vulnerabilities accumulated for four decades because
discipline fails at scale.
Rust's ownership model solved the same problem by making the
\emph{absence of a deallocation owner} a compile-time error, not a runtime check.
The lesson generalizes:
no existing framework makes the absence of evidence a compile-time error
rather than a runtime omission.

Linear Logic, introduced by Girard in 1987~\cite{girard1987linear} and
connected to programming language semantics by Wadler~\cite{wadler1991linear},
treats propositions as \emph{resources}: they must be used exactly once,
can be neither duplicated nor discarded.
The connective $\multimap$ (linear implication) denotes a resource transformation
that \emph{consumes} its premise, and $\otimes$ (tensor product) denotes the
\emph{joint} consumption of two resources---not their independent conjunction.
We exploit this semantics to encode the accountability invariant directly
in the type system.
Define the type law $V \multimap (E \otimes C)$:
a Verdict~($V$) can only be produced by linearly consuming an
\texttt{EvidenceToken}~($E$) and a \texttt{ComplianceToken}~($C$) together.
In our Rust implementation, \texttt{EvidenceToken} is a \texttt{\#[must\_use]}
struct with no \texttt{Clone} or \texttt{Copy} derivation;
\texttt{Verdict::new()} takes both tokens by value, consuming them atomically.
Because \texttt{Blake3Hash}---the only evidence artifact---has no public
constructor, it cannot be forged: the sole path to a \texttt{Blake3Hash}
is through \texttt{EvidenceToken::consume()}.
A~silent decision is therefore not an illegal act that slips past a guard;
it is a \emph{type error, not a runtime omission}:
the compiler rejects it before a single byte of evidence is ever dropped.

This paper makes three contributions.
\begin{enumerate}
  \item \textbf{Linear Resource Types for AI Accountability (C1).}
    We present the first formal encoding of AI accountability obligations---
    specifically LGPD Art.~18\S2 and EU AI Act Art.~86---as linear resources
    in the sense of Girard~\cite{girard1987linear}.
    The core invariant $V \multimap (E \otimes C)$ is expressed as a
    compile-time type constraint in Rust and verified on the companion
    repository~\cite{soares2026a, soares2026b}.

  \item \textbf{Decision-Evidence Causal Binding (C2).}
    We prove that, under the \emph{BTV}\footnote{BTV (\textit{Bound-To-Verdict}):
    a reference kernel implementing atomic accountability via linear types.
    Evidence and compliance tokens are structurally \emph{bound to} the
    verdict --- the name is the type law $V \multimap (E \otimes C)$ in English.}
    type invariant, every materialized
    \texttt{Verdict} consumed \emph{exactly one} causally-prior
    \texttt{EvidenceToken}: a BLAKE3 digest of the decision context
    computed before the verdict existed.
    Linear ownership prevents duplication; the private constructor
    prevents forgery; both guarantees are static.

  \item \textbf{The Constitutional Enclosure Theorem (C5).}
    We state and prove that in any system satisfying the BTV type invariant,
    the set of materialized silent decisions is empty.
    The proof proceeds by contradiction via the Rust type system and is
    machine-checkable: \texttt{cargo test} on the companion repository
    verifies all six clauses in under four seconds.\
\end{enumerate}

\noindent
Section~2 reviews the necessary background on Linear Logic and Rust ownership.
Section~3 defines the BTV type system formally.
Section~4 states and proves the Constitutional Enclosure Theorem.
Section~5 describes the implementation and benchmarks.
Section~6 discusses related work and limitations.
